\section{States with Biggest Improvement Over Enacted Maps}
\label{sec:enacted}

The comparison of algorithmic maps to enacted maps answers a different question than the scale and resolution comparisons: not ``how compact can algorithmic redistricting be?'' but ``how much better than what legislators actually draw?'' We compile enacted state house district boundaries from publicly available GIS sources for 35 states where current or recent enacted maps are available in TIGER-compatible format, and compute Polsby-Popper ratios for both the enacted and algorithmic maps.

\subsection{Data Sources}

Enacted state legislative district shapefiles were obtained from the following sources: state legislative GIS offices (for 22 states), the All About Redistricting database maintained by Justin Levitt (for 8 states), and the Redistricting Data Hub (for 5 states). All shapefiles were projected to the Albers Equal Area Conic projection before computing PP ratios to ensure area and perimeter calculations are not distorted by geographic projection. The 15 states for which current enacted maps were not available in machine-readable GIS format are excluded from this comparison.

\subsection{Mean Improvement Across All States}

Across the 35 states with available data, the algorithmic state house maps (block group resolution) outperform the enacted state house maps by a mean of 0.062 PP units, which represents a 17\% improvement over the enacted maps' mean PP of 0.362. The comparison is:

\begin{itemize}
\item \textbf{Enacted state house maps}: mean PP = 0.362 (SD: 0.051)
\item \textbf{Algorithmic state house maps (BG)}: mean PP = 0.401 (SD: 0.037)
\item \textbf{Improvement}: +0.039 mean, +17\%, statistically significant ($t(34) = 6.1$, $p < 0.001$)
\end{itemize}

The algorithmic maps outperform enacted maps in 31 of 35 states. The four states where enacted maps have higher PP than algorithmic maps are all small states (Rhode Island, Vermont, Delaware, South Dakota) where the enacted maps are likely drawn by legislators with knowledge of local geographic features that the algorithm cannot access.

\subsection{States with Largest Algorithmic Improvement}

The states with the largest algorithmic improvement over enacted maps (sorted by improvement in PP units) are states where the enacted maps have been identified as gerrymandered or where legislative redistricting is highly politicized:

\begin{table}[htbp]
\centering
\caption{States with largest algorithmic improvement over enacted state house maps.}
\label{tab:enacted_improvement}
\begin{tabular}{lrrrr}
\toprule
\textbf{State} & \textbf{Enacted PP} & \textbf{Algorithmic PP} & \textbf{Improvement} & \textbf{Notes} \\
\midrule
Maryland    & 0.201 & 0.334 & +0.133 & Court-identified gerrymander \\
North Carolina & 0.248 & 0.379 & +0.131 & Harper v. Hall (2022) \\
Pennsylvania & 0.253 & 0.387 & +0.134 & LWV struck down 2018 \\
Ohio        & 0.271 & 0.384 & +0.113 & Amendment 1 violations\\
Illinois    & 0.282 & 0.387 & +0.105 & Legislative gerrymander \\
Texas       & 0.296 & 0.371 & +0.075 & VRA litigation ongoing \\
Wisconsin   & 0.318 & 0.408 & +0.090 & Court-ordered redraw \\
New York    & 0.298 & 0.368 & +0.070 & IRC established post-litigation \\
Michigan    & 0.324 & 0.389 & +0.065 & Pre-MICRC enacted map \\
Florida     & 0.329 & 0.393 & +0.064 & Amendment 6 violations \\
\midrule
\textbf{Mean (top 10)} & 0.282 & 0.380 & +0.098 & \\
\textbf{Mean (all 35)} & 0.362 & 0.401 & +0.039 & \\
\bottomrule
\end{tabular}
\end{table}

\noindent Table~\ref{tab:enacted_improvement} reveals a clear pattern: the states with the largest algorithmic improvement are disproportionately states that courts have identified as having gerrymandered maps. Pennsylvania's 2011 state house map was struck down by the state supreme court in \textit{League of Women Voters v.\ Commonwealth}; its enacted PP of 0.253 is among the lowest of any state in the dataset. North Carolina's maps were struck down in \textit{Harper v.\ Hall} on partisan gerrymandering grounds; the enacted PP of 0.248 reflects the extreme non-compactness of the gerrymandered districts.

This correlation between low enacted PP and identified gerrymandering is consistent with the argument that compactness is a proxy for partisan manipulation \citep{polsby1991third}. The extreme non-compactness of gerrymandered maps---Maryland's 0.201, Pennsylvania's 0.253, North Carolina's 0.248---reflects the tortured boundary drawing required to pack or crack minority partisan communities. The algorithmic maps, which cannot see partisan data, produce substantially more compact configurations.

\subsection{States Where Enacted Maps Are Competitive}

In 4 states, enacted maps outperform algorithmic maps on compactness:

\begin{itemize}
\item \textbf{Rhode Island}: Enacted PP 0.481 vs.\ algorithmic 0.474. Rhode Island's small size and two-district congressional map mean that the legislature has few opportunities to gerrymander, and local geographic knowledge may produce slightly more compact districts.
\item \textbf{Vermont}: Enacted PP 0.505 vs.\ algorithmic 0.492. Vermont's single congressional district means all variation is in state house; the legislature's granular local knowledge produces slightly higher PP.
\item \textbf{Delaware}: Enacted PP 0.512 vs.\ algorithmic 0.503. Same reasoning as Vermont; single at-large congressional district.
\item \textbf{South Dakota}: Enacted PP 0.461 vs.\ algorithmic 0.449. Plains geography with regular tract shapes; local knowledge advantage is not offset by any gerrymandering disadvantage.
\end{itemize}

These results are consistent with the finding in \citet{altman2011promise} that algorithmic redistricting's advantage is most pronounced in states with more districts and more political competition. In small states with limited gerrymandering opportunities, the algorithm's advantages are smaller and local knowledge may provide some offset.
